% =====================================================================
%  GLOSARIO Y ABREVIATURAS
%
%  POR QUE ESTE GLOSARIO NO REPITE EL MARCO CONCEPTUAL
%  ---------------------------------------------------
%  El §2.2 define cincuenta términos agrupados por tema, y es el lugar donde se
%  explican a fondo. Un glosario que los copiara sería redundante, y la
%  redundancia ya se había señalado como problema del documento.
%
%  Este glosario tiene otra función: permitir leer el resumen y la introducción
%  SIN haber llegado al capítulo 2. Por eso trae solo los términos que aparecen
%  antes de ese capítulo, en orden alfabético, con la definición mínima que hace
%  falta, y remite al §2.2 para el desarrollo.
%
%  La plantilla sugiere ubicarlo como anexo. Se deja en los preliminares porque
%  es donde lo pone su propia tabla de contenido, y porque un glosario que se
%  lee después del documento no cumple su función.
% =====================================================================
\ustaseccion{Glosario}

Los términos siguientes se usan desde el resumen y la introducción. El
capítulo~2 los desarrolla y añade los demás.

\begin{description}[style=nextline,leftmargin=1.5em,itemsep=0.2em]
  \item[Agente] Cada uno de los dos robots del sistema. Se usa «agente» y no
    «robot» cuando importa su papel dentro de la coordinación y no su hardware.

  \item[AMCL] Algoritmo de localización por filtro de partículas. Combina la
    odometría con las lecturas del láser para situar al robot dentro de un mapa
    que ya conoce. No construye el mapa: lo supone dado.

  \item[Banda muerta de tracción] Rango de velocidades ordenadas en el que el
    vehículo no llega a moverse, porque la señal resultante no vence el
    rozamiento estático. En la plataforma de este trabajo abarca todo lo que
    está por debajo de 0,40~m/s, y condiciona la aproximación final del vehículo
    al destino.

  \item[Cinemática tipo automóvil (dirección Ackermann)] Modelo de un vehículo
    que dirige con el eje delantero. Tiene radio de giro mínimo y no puede girar
    sobre su propio eje. Es lo contrario de la tracción diferencial, que es lo
    que la mayoría de las herramientas de navegación supone por omisión.

  \item[Continuidad del servicio] Propiedad de no interrumpir la asistencia
    mientras el usuario cambia de nivel. Es el objeto de la pregunta problema.

  \item[Localización] Estimación de la posición y la orientación del robot
    respecto a un marco de referencia fijo.

  \item[Mapa de ocupación] Cuadrícula en la que cada celda se declara libre,
    ocupada o desconocida. Es la representación del edificio con la que
    funcionan la localización y la planificación.

  \item[Mapa de costos (\textit{costmap})] Cuadrícula derivada de la anterior en
    la que cada celda lleva un costo de tránsito en lugar de una sola etiqueta.
    Permite que el planificador prefiera el centro del pasillo sin prohibir
    los bordes.

  \item[Middleware robótico] Capa de software que comunica entre sí los módulos
    del sistema. Aquí es ROS~2.

  \item[Misión] Una solicitud de guiado completa, desde que el usuario elige
    origen y destino hasta que el sistema la cierra. Es la unidad de medida de
    la campaña experimental: las métricas se reportan por misión.

  \item[Observabilidad] Propiedad de una magnitud física que permite estimarla a
    partir de las mediciones disponibles. Si en cierta combinación de entorno y
    sensor una magnitud es inobservable, ningún ajuste de parámetros la
    recupera, porque la información no está en la señal. El §2.3.4 muestra un
    caso concreto en un pasillo recto.

  \item[Odometría] Estimación del desplazamiento propio del robot a partir de
    sus propias mediciones. La plataforma de este trabajo no tiene encoders, así
    que la obtiene comparando barridos sucesivos del láser.

  \item[Punto de transferencia] Lugar físico, en este trabajo la escalera, donde
    un agente entrega el guiado y el otro lo recibe.

  \item[Relevo] Traspaso de la responsabilidad de guiar a una persona, de un
    agente al otro, en un punto acordado de antemano. Es el aporte que el
    proyecto declara, y lo que sustituye al hardware de transición vertical.

  \item[Tasa de éxito] Proporción de misiones que terminan como el protocolo
    experimental declara que deben terminar. Su definición operativa está en el
    §3.4.

  \item[Tiempo de respuesta] Intervalo entre la solicitud del usuario y el
    instante en que el robot asignado arranca la tarea.

  \item[Zona de transición] Punto del edificio donde se cambia de nivel:
    escalera o ascensor. En este sistema ningún robot la atraviesa.
\end{description}

\ustaseccion{Abreviaturas}

\begin{description}[style=nextline,leftmargin=1.5em,itemsep=0.2em]
  \item[AMCL] \textit{Adaptive Monte Carlo Localization} — localización adaptativa
    por método de Monte Carlo.
  \item[AWS] \textit{Amazon Web Services}, fabricante de la plataforma DeepRacer
    empleada como vehículo.
  \item[CONPES] Consejo Nacional de Política Económica y Social (Colombia).
  \item[DDS] \textit{Data Distribution Service} — protocolo de transporte sobre
    el que ROS~2 implementa sus comunicaciones.
  \item[GED] Grupo de Estudio y Desarrollo en Robótica, grupo de investigación en
    el que se inscribe el proyecto.
  \item[GUI] \textit{Graphical User Interface} — interfaz gráfica de usuario.
  \item[HRI] \textit{Human-Robot Interaction} — interacción humano-robot.
  \item[ICONTEC] Instituto Colombiano de Normas Técnicas y Certificación.
  \item[IMU] \textit{Inertial Measurement Unit} — unidad de medición inercial.
  \item[JSON] \textit{JavaScript Object Notation}, formato de intercambio de
    datos en texto.
  \item[MRS] \textit{Multi-Robot System} — sistema multi-robot.
  \item[Nav2] \textit{Navigation2}, marco de navegación autónoma de ROS~2.
  \item[NTC] Norma Técnica Colombiana.
  \item[OE] Objetivo específico. Los cuatro del proyecto están en el §1.3.2 y
    organizan los capítulos~3 y~4.
  \item[PID] Control proporcional-integral-derivativo.
  \item[PWM] \textit{Pulse Width Modulation} — modulación por ancho de pulso.
  \item[QoS] \textit{Quality of Service} — políticas de entrega de un canal de
    comunicación.
  \item[RF] Requisito funcional. La numeración RF-\textit{nn} remite a la
    especificación de requisitos del sistema.
  \item[ROS] \textit{Robot Operating System}. La versión empleada es ROS~2, que
    no es una actualización de ROS~1 sino un rediseño sobre DDS.
  \item[SLAM] \textit{Simultaneous Localization and Mapping} — localización y
    mapeo simultáneos.
\end{description}
